\documentclass{article}
\usepackage[UTF8]{ctex}
\usepackage{amsmath,amssymb,amsfonts}
\usepackage{enumitem}
\usepackage{geometry}
\geometry{a4paper, margin=1in}
% === 错题本排版优化（不参与题目标识） ===
\setlength{\parindent}{0pt}
\setlength{\parskip}{0.35em plus 0.10em minus 0.08em}
\setlength{\emergencystretch}{3em}
\setlist{itemsep=0.25em, topsep=0.40em, parsep=0pt}
\allowdisplaybreaks[2]
% === 错题本排版优化结束 ===
\begin{document}
\section*{南京理工大学2024年数学分析真题}
\section*{96．南京理工大学2024年数学分析试卷}

\subsection{题目1}
一．（15 分）解答如下问题：

\begin{itemize}
  \item[（1）] 求极限 \(\lim _{x \rightarrow \infty}\left(\cos \frac{1}{x}\right)^{x^{2}-1}\) ．
\end{itemize}

\subsection{题目5}
原题边界“（2）”：设 \(z=x^{2}+y^{2}\) ，其中 \(y=f(x)\) 是由曲线 \(x^{2}-x y+y^{2}=1\) 确定的隐函数，求 \(\frac{\mathrm{d} z}{\mathrm{~d} x}\) ．

\subsection{题目2}
二．（15 分）解答如下问题：

\begin{itemize}
  \item[（1）] 判断 \(\sum_{n=1}^{\infty} \frac{1}{\left[\ln \left(\tan \frac{1}{n}\right)\right]^{3}}\) 的敛散性．
\end{itemize}

\subsection{题目6}
原题边界“（2）”：判断 \(\int_{1}^{+\infty} \frac{\sin x^{2}}{1+x^{p}} \mathrm{~d} x(p \geq 0)\) 的敛散性．

\subsection{题目3}
三．（15 分）设 \(f(x)\) 在 \(\mathbb{R}\) 上二阶可导，且 \(f^{\prime}(0)=0\) ，证明：存在 \(\xi \in(0,1)\) ，使得 \(f^{\prime}(\xi)=(\xi-1)^{2} f^{\prime \prime}(\xi)\) ．

\subsection{题目4}
四．（15 分）设 \(f(x)=\sum_{n=1}^{\infty} n e^{-n x}\) ．

\begin{itemize}
  \item[（1）] 讨论 \(f(x)\) 的连续区间．
  \item[（2）] 计算 \(\int_{1}^{2} f(x) \mathrm{d} x\) ．
\end{itemize}

\subsection{题目5}
五。（15 分）设 \(f(x)\) 在 \(\mathbb{R}\) 上连续，且 \(\lim _{x \rightarrow \infty}[f(x)-\lambda x]=0\) ，其中 \(\lambda\) 为常数，证明：\(f(x)\) 在 \(\mathbb{R}\) 上一致连续。

\subsection{题目6}
六。（15分）解答如下问题：

\begin{itemize}
  \item[（1）] 求 \(f(x)=\arctan \frac{x+3}{x-3}\) 的麦克劳林级数展开式，并讨论收敛区间．
\end{itemize}

\subsection{题目6}
原题边界“（2）”：求 \(f(x)=x^{3}\) 在 \([0, \pi]\) 上的余弦级数．

\subsection{题目7}
七．（15 分）计算三重积分

\[
\iiint_{V} z^{2} \mathrm{~d} x \mathrm{~d} y \mathrm{~d} z
\]

其中 \(V: x^{2}+y^{2}+z^{2} \leq a^{2}, x^{2}+y^{2} \leq a x, a>0\) ．

\subsection{题目8}
八．（15 分）计算曲面积分

\[
I=\iint_{S} x^{2} \mathrm{~d} y \mathrm{~d} z+y^{2} \mathrm{~d} z \mathrm{~d} x+z^{2} \mathrm{~d} x \mathrm{~d} y .
\]

其中 \(S\) 为 \(x^{2}+y^{2}=z^{2}, 0 \leq z \leq 1\) ，取曲面下侧。

\subsection{题目9}
九．（15 分）计算曲线积分

\[
\lim _{d(\Omega) \rightarrow 0} \frac{1}{S(\Omega)} \oint_{L}(\mathbf{F}, \mathbf{n}) \mathrm{d} s .
\]

其中 \(\Omega\) 是包含原点且由简单封闭的光滑曲线 \(C\) 围成的区域，\(d(\Omega)\) 为 \(\Omega\) 的直径，\(S(\Omega)\) 为 \(\Omega\) 的面积\(\mathbf{F}=(P(x, y), Q(x, y))\) 为区域 \(\Omega+C\) 上连续可微的向量函数， \(\mathbf{n}\) 为曲线 \(C\) 的单位外法向量．

\footnotetext{十．（15 分）用有限覆盖定理证明致密性定理，即任意有界数列 \(\left\{x_{n}\right\}\) 必定有收敛子列．
}
\end{document}
